% !TeX root = ../main.tex
\chapter{Introduction}
\label{chap:introduction}

\textbf{Draft outline.} Introduce the problem, motivation, research question,
objectives, contributions, and organization of the thesis.

% Answer: what problem am I adressing and why is it important?

% Existing material: ../plano-de-trabalho/sections/01-introducao.tex
% and ../tcc-docs/README.md. Translate/adapt deliberately as the thesis develops.

% First, talk about the problem of having to iterate fast on dynamic and agile environnemnts.

%

% TODO: Write the motivation and problem before this section; write the
% contributions and thesis organization after Chapters 3 and 6 are drafted.

\section{Research Questions}
\label{sec:research-questions}

This thesis addresses the following research question:

\begin{quote}
How can a pattern that connects evaluation, monitoring, and curation support
frequent delivery of AI-enabled systems while preserving established behavior?
\end{quote}

The pattern examined is CD4AI, presented in Chapter~\ref{chap:cd4ai}. The
question is divided into three sub-questions:

\begin{description}
  \item[RQ1.] How do the pattern's stages and their connections appear in an
    implementation and in existing tools?
  \item[RQ2.] How do these cases connect evaluation results to delivery
    decisions, in particular when distinguishing known limitations from
    regressions?
  \item[RQ3.] What do the cases suggest for refining the pattern?
\end{description}

% Where each question is answered (for writing; not in the final text):
% RQ1 and RQ2: Chapter 5 (feedback loop; result interpretation and delivery).
% RQ3: Chapter 6 (implications for CD4AI). Summary of answers: Chapter 7.
